% 2026-11-12 / B-v44: first-order cancellation, independently of the all-orders transport.
\section{一次の非可換変形と可換模型の接空間}
\label{fpk:sec:first-order}

この節では一般点の全次数形式比較を前提とせず、
双対数環上で非可換積の一次摂動が曲率付きMFにどう作用するかを直接調べる。
これは全次数の定理\ref{cft:thm:complete}の一次部分の独立検証でもある。
曲率の一次変形を固定しなければ、任意の非可換変形が自明という意味ではない。

$S=\mathbb C[x_1,\ldots,x_N]$、$N\ge5$、
$D=\mathbb C[\varepsilon]/(\varepsilon^2)$ とする。
$F\in S_3$は滑らかな射影三次超曲面を定めるものとする。
$A_D$を単位付き・次数別有限自由な$D$-代数で
$A_D/\varepsilon A_D\simeq S$、
$\Omega_D\in Z(A_D)_3$の剰余が$F$となるものとする。

\begin{theorem}[一次Poisson方向の曲率付き消滅]
\label{fpk:thm:infinitesimal}
上の入力に対し、単位と内部次数を保つ$D$-加群標識を適切に選べば、
$A_D$の積と中心元は
\begin{equation}\label{fpk:eq:normal}
 a*b=ab+\varepsilon\,\Theta_\pi(a,b),\qquad
 \Omega_D=F+\varepsilon G,\qquad
 \Theta_\pi(a,b)=\frac12\{a,b\}_\pi
\end{equation}
と書ける。ここで$\pi$は二次係数のbivectorで$[\pi,F]=0$である。
さらに、この標識で有限rank graded MFの
$D$-線形unital $A_\infty$同値
\begin{equation}\label{fpk:eq:equivalence}
 \operatorname{MF}^{\mathbb Z}(A_D,\Omega_D)
 \simeq
 \operatorname{MF}^{\mathbb Z}(S\otimes D,F+\varepsilon G)
\end{equation}
が存在し、剰余では恒等、内部shiftおよび
中心元による三次周期同型$(3)\simeq[2]$を保つ。

二つの標識で得た$G$の差は
$J(F)_3=(\partial_1F,\ldots,\partial_NF)_3$に属する。
従って、一次の可換模型の座標同値類を定める写像は
\begin{equation}\label{fpk:eq:tangent}
 [G]\in(S/J(F))_3
\end{equation}
を通る。特に$G=0$と取れる正規化では、一次Poisson変形は
rotation付きgraded MF圏の一次変形を生じない。
\end{theorem}

\begin{proof}
次数別自由な加群標識を選ぶと、積は$\mu+\varepsilon\beta$となる。
単位を保つよう標識を取り、$\beta$をnormalizedな内部weight零の
Hochschild二cochainとする。結合律の一次部分は
$d_{\mathrm{Hoch}}\beta=0$である。
polynomial algebraのHKR同型により、そのクラスは
二次係数bivector$\pi$の反対称微分cochain$\Theta_\pi$で表せる。
差$\beta-\Theta_\pi$は内部weight零のHochschild coboundaryなので、
$\mathrm{id}+\varepsilon T$という一次の厳密な加群標識変更で
積を式\eqref{fpk:eq:normal}にできる。
この変更は元の代数の結合律を保ち、単位も保つ。
中心元の表示を$F+\varepsilon G$と書く。
中心性の一次部分は$[\pi,F]=0$である。

$F$が滑らかなので、全偏微分は正則列である。
$\iota_{dF}$による斉次Koszul複体をrank二・係数次数二で用いると、
$[\pi,F]=0$から定数係数三重vector
$\Xi\in\bigwedge^3\langle\partial_1,\ldots,\partial_N\rangle$を得て
\begin{equation}\label{fpk:eq:primitive}
 \pi=[\Xi,F]
\end{equation}
となる。原像の差を消すにはrank四・係数次数$-2$の原像が必要となり、
そのような多項式はないので$\Xi$は一意である。
ここでは$\pi$のJacobi恒等式を使わない。

$\Theta_\Xi=\operatorname{HKR}(\Xi)$を、標準の$1/3!$反対称化で
正規化した三入力Hochschild cochainとする。
三つの座標方向の単純三重vectorについては
\[
 \Theta_\Xi(a,b,c)
 =\frac16\det
 \begin{pmatrix}
  \partial_i a&\partial_i b&\partial_i c\\
  \partial_j a&\partial_j b&\partial_j c\\
  \partial_k a&\partial_k b&\partial_k c
 \end{pmatrix}.
\]
式\eqref{fpk:eq:primitive}と反対称性より
\begin{equation}\label{fpk:eq:insert}
 \Theta_\Xi(F,a,b)-\Theta_\Xi(a,F,b)+\Theta_\Xi(a,b,F)
 =\Theta_\pi(a,b).
\end{equation}
$\Theta_\Xi$は多重導分なので
$d_{\mathrm{Hoch}}\Theta_\Xi=0$である。

周期変数に$\deg_{\mathrm{coh}}t=2$、
$\deg_{\mathrm{int}}t=-3$を与え、曲率を$-tF$とする。
\[
 \eta=\varepsilon\,t^{-1}\Theta_\Xi
\]
はcohomological・内部次数ともに零、arity三のnormalized cochainである。
曲率付きGerstenhaber括弧で式\eqref{fpk:eq:insert}を使うと
\[
 [\eta,\mu]=0,\qquad
 [\eta,-tF]=-\varepsilon\Theta_\pi
\]
となる（符号は式\eqref{fpk:eq:insert}を採用する）。
従って、曲率付き構造
$\mu-tF+\varepsilon(\Theta_\pi-tG)$は、
$\eta$による一次ゲージで$\mu-t(F+\varepsilon G)$へ移る。
$\varepsilon^2=0$だから、そのbar写像は
$F_0=0,F_1=\mathrm{id},F_3=\varepsilon t^{-1}\Theta_\Xi$のみを持ち、
逆は$F_3$の符号を反転させればよい。

有限rank MFの行列$Q$は$\varepsilon$の倍数である必要はない。
しかし対象像$Q'=Q+F_3(Q,Q,Q)$および射のtwisting挿入公式は
双対数環上では有限和である。
bar恒等式がMF方程式と関手恒等式を与え、逆bar写像で互いに逆な
$A_\infty$関手が得られる。写像は内部weight零で
$t$を微分しないので内部shiftと$(3)\simeq[2]$を保つ。
これで式\eqref{fpk:eq:equivalence}が従う。

最後に、積を$\mu+\varepsilon\Theta_\pi$へ正規化する
二つの標識の差を$\mathrm{id}+\varepsilon T$とする。
同じ正規化された積を保つ条件は
$d_{\mathrm{Hoch}}T=0$であるから、$T$は内部weight零の導分である。
従って$T(x_i)=\sum_j a_{ij}x_j$であり、
中心元の一次係数は$G\mapsto G+T(F)$と変わる。
この差は$J(F)_3$に属する。
また、同じ$F$の一次座標変更による変形は
この形の$T(F)$で尽くされる。
従って式\eqref{fpk:eq:tangent}を得る。
\end{proof}

\begin{corollary}[Jacobian接空間の次元]
\label{fpk:cor:dimension}
滑らかな三次式$F$に対して
\begin{equation}\label{fpk:eq:dimension}
 \dim_{\mathbb C}(S/J(F))_3
 =\binom{N+2}{3}-N^2.
\end{equation}
特にcubic threefoldでは$10$、cubic fourfoldでは$20$である。
\end{corollary}
\begin{proof}
偏微分の二次式は正則列であり、一次係数のsyzygyを持たない。
従って$J(F)_3$は$N$個の二次偏微分に$N$次元の一次式を掛けた
$N^2$次元の部分空間である。
一方$\dim S_3=\binom{N+2}{3}$である。
\end{proof}

\begin{remark}[Jacobi条件が必要になる位置]
\label{fpk:rem:jacobi}
一次の曲率付き同値に必要なのは中心性からの$[\pi,F]=0$であり、
$[\pi,\pi]=0$は必要ない。
二次以降で同じゲージを延長するには
$[\Xi,\pi]=0$が必要になる。
滑らかな三次式の下では、Jacobi条件$[\pi,\pi]=0$から
勾配Koszulと係数次数により$[\Xi,\pi]=0$を導ける
（補題\ref{cft:lem:primitive}）。
従って一次消滅と全次数の比較は異なる論理的強さを持つ。
\end{remark}


\subsection{一次対変形の接空間と二次のPl\"ucker障害}

\begin{proposition}[次数付き中心三次対の一次接空間]
\label{fpk:prop:tangent-pair}
$N\ge5$とし、可換対$(S,F)$における次数別flatな中心三次対の
双対数変形を、剰余で恒等となる単位付き次数付き同型で割る。
中心元は次数三の実際の元として標識する。
この一次変形集合には自然なベクトル空間構造があり、
HKR正規形を選べば
\begin{equation}\label{fpk:eq:tangent-pair}
 T^1_{\mathrm{pair},(S,F)}
 \simeq
 \bigwedge^3\mathbb C^N\oplus\operatorname{Jac}(F)_3
\end{equation}
と表される。第一成分は$\pi=\iota_{dF}\Xi$、第二成分は
$\Omega=F+\varepsilon G$の$[G]$である。
特に$N=5$ではこの接空間は$20$次元である。
\end{proposition}
\begin{proof}
定理\ref{fpk:thm:infinitesimal}で用いたnormalized HKRにより、
積の一次項の同型類は二次係数bivector$\pi$で記述される。
中心性の一次条件は$[\pi,F]=0$であり、勾配Koszulの完全性と
係数次数から唯一の定数三重vector$\Xi$が存在して$\pi=[\Xi,F]$となる。
逆に任意の$\Xi$はこの中心性条件を満たし、
$\Theta_\pi$を一次積とする双対数上の結合代数を与える。
任意の三次式$G$は中心元の一次項に取れる。
一方、積のHKR正規形を保つ剰余恒等の次数付き標識変更は
一次の導分$T$であり、$G$を$G+T(F)$へ移す。
内部次数零の導分は一次形式の線形変換だから、
$T(F)$の張る空間は$J(F)_3$である。
よって表示を得る。
\end{proof}

以下では$N=5$とする。固定体積形式を使い、
$\bigwedge^3V\simeq\bigwedge^2V^*\otimes\det V$
で$\Xi$に対応する二形式を$\alpha_\Xi$と書く。
その五つのPfaffianを$q_0,\ldots,q_4$とし、
\[
 q(\Xi)=\alpha_\Xi\wedge\alpha_\Xi
 \in\bigwedge^4V^*
\]
で表す。$q(\Xi)=0$は$\Xi$の分解可能性と同値で、
$q_i$は$\operatorname{Gr}(3,V)$のPl\"uckerイデアルを生成する。

\begin{theorem}[二次Jacobi障害のスキーム論的Pl\"ucker表示]
\label{fpk:thm:plucker-obstruction}
$V=\mathbb C^5$、$F\in\operatorname{Sym}^3V^*$を
滑らかな射影三次超曲面の方程式とする。
定数三重vector$\Xi$について$\pi_\Xi=\iota_{dF}\Xi$と置くと
\[
 [\pi_\Xi,\pi_\Xi]=0
 \quad\Longleftrightarrow\quad
 \Xi\text{は分解可能（零を含む）}.
\]
さらに$[\pi_\Xi,\pi_\Xi]$の全係数を$\Xi$の座標の二次式と見て
生成されるイデアルは、五つのPl\"ucker二次式のイデアル
$(q_0,\ldots,q_4)$に\emph{スキームとして}一致する。
\end{theorem}
\begin{proof}
まず$\Xi$が分解可能なら、座標変換により
$\Xi=\partial_0\wedge\partial_1\wedge\partial_2$とできる。
$\pi_\Xi$は三つのactive変数のJacobian bracketであり、
任意の$F$についてJacobi恒等式を満たす。

逆に$\Xi$が分解不可能なら、五次元における二形式の交代標準形より、
座標変換と非零スカラー倍で
\[
 \Xi=\partial_0\wedge\partial_1\wedge\partial_2+
       \partial_0\wedge\partial_3\wedge\partial_4
\]
とできる。$F_i=\partial_iF$、$F_{0i}=\partial_0\partial_iF$
と書く。$\pi_{ij}=\sum_k\Xi_{ijk}F_k$を用いてJacobiatorを直接展開すると、
$[\pi_\Xi,\pi_\Xi]=0$はちょうど
\begin{equation}\label{fpk:eq:jacobi-normal}
 F_{0i}F_j-F_{0j}F_i=0
 \qquad(0\le i<j\le4)
\end{equation}
に等しい（成分の順序による符号と共通非零係数は無関係）。
式\eqref{fpk:eq:jacobi-normal}に$x_j$を掛けて$j$について加え、
Euler等式$\sum_jx_jF_j=3F$および
$\sum_jx_jF_{0j}=2F_0$を使えば
\begin{equation}\label{fpk:eq:euler-obstruction}
 3F F_{0i}=2F_0F_i
 \qquad(0\le i\le4)
\end{equation}
を得る。もし$F_0=0$なら$F$は$x_0$に依らず、
$[1:0:\cdots:0]$が特異点となる。
$F_0\ne0$なら、滑らかな射影三次式$F$は既約なので
$\gcd(F,F_0)=1$である。式\eqref{fpk:eq:euler-obstruction}から
$F\mid F_i$が全$i$について従うが、$\deg F_i=2<3$なので
全$F_i=0$となり矛盾する。従って$\Xi$は分解可能である。

次に二次式の\emph{イデアル}の一致を示す。
$[\pi_\Xi,\pi_\Xi]$の係数は$\Xi$について斉次二次式であり、
分解可能な$\Xi$で零になる。Grassmannianの斉次イデアルは
五つのPl\"ucker二次式で生成されるので、
係数全体は$q_0,\ldots,q_4$の線形結合である。
従ってある線形写像$L_F$について、Jacobiatorの全係数ベクトルは
$L_F(q(\Xi))$と表される。

ここで二次写像$q:\bigwedge^3V\to\bigwedge^4V^*$は全射である。
実際、五次元では任意の非零四形式は分解可能であり、
適当な基底で$c\,e^1\wedge e^2\wedge e^3\wedge e^4$と書ける。
$\alpha=e^1\wedge e^2+(c/2)e^3\wedge e^4$とすれば
$\alpha\wedge\alpha$はその四形式である。
もし$L_F$に非零の核があれば、$q(\Xi)$がその非零元となる$\Xi$を選べる。
すると$[\pi_\Xi,\pi_\Xi]=0$なのに$q(\Xi)\ne0$であり、
前半の結論と矛盾する。よって$L_F$は単射である。
その双対写像の全射性から、Jacobiatorの係数二次式は
Pl\"ucker二次式の五次元空間全体を張り、生成イデアルも一致する。
\end{proof}

\begin{corollary}[二次まで延長できる非可換一次方向]
\label{fpk:cor:second-order}
$N=5$の中心三次対の一次変形の非可換方向
$\Xi\in\bigwedge^3V$が、次数付き\emph{結合積}の二次変形
$\mathbb C[\varepsilon]/(\varepsilon^3)$へ延長可能であるためには、
五つのPl\"ucker二次式が零でなければならない。
結合積に関する二次障害の零スキームはPl\"ucker錐と一致する。
その錐の局所次元は$7$で、頂点のZariski接空間は$10$次元である。
これは中心元を含めた対の高次延長障害がすべて消えるという主張ではない。
\end{corollary}
\begin{proof}
結合積$\mu+\varepsilon\beta_1+\varepsilon^2\beta_2$の結合律の
二次項は
$d_{\rm Hoch}\beta_2+\frac12[\beta_1,\beta_1]=0$である。
HKR同型で二次障害の$HH^3$類は
$\frac12[\pi_\Xi,\pi_\Xi]$に対応し、
この類が零であることが結合積の二次延長の必要十分条件である。
前定理からその係数イデアルは五つのPl\"ucker関係式になる。
$\operatorname{Gr}(3,5)$の次元は$3(5-3)=6$なので
錐の次元は$7$、零点では全五式の一次微分が零で接空間は$10$次元である。
\end{proof}

\begin{remark}[既知結果と本研究への意味]
\label{fpk:rem:plucker-scope}
滑らかなcubic threefoldのPoisson構造のGrassmannian表示は
既知の研究\cite{PMayanskiy}と整合する。
ここで新たに整理したのは、一次中心三次対の変形から二次Hochschild障害へ
進む際に、この五つの二次式が\emph{イデアルとして過不足なく}現れる
直接的な証明である。この説明自体の初出は未確認である。
一次接空間の次元$20$と、二次障害零スキームの局所次元$17$を混同しない。
後者の全高次の形式局所環との一致や、非零Poisson点での
categoricalな忘却射の核の同定は、別途検証する必要がある。
\end{remark}

\begin{remark}[17次元族と可換moduliの関係]
\label{fpk:rem:moduli}
$cubic\ threefold$の次数三Jacobian空間の次元は$10$である。
この事実は、非可換モジュライの余分な一次方向が
rotation付き残余圏では不可視になり得ることを説明する。
ここで、emph{一般点の次元}とemph{可換点のZariski接空間}を区別する。
固定した滑らかな五変数三次式に対するPoisson錐をPl"ucker錐と同定すれば、
非零分解可能な一般点でその次元は
$1+dimoperatorname{Gr}(2,5)=7$である。
一方、その頂点ではPl"ucker方程式は二次式なのでZariski接空間は
$igwedge^3mathbb C^5$の全$10$次元となる。
従って、可換三次式の接空間$10$次元と合わせた形式的対の
接空間は$20$次元であり、局所的な成分次元$17$と一致する必要はない。
	exttt{nc-cubic2} の中心三次対の全Artin局所計算も
「局所次元$17$、接空間次元$20$」を区別している。

本定理が消すのはemph{可換点における}これらの
一次非可換方向である。一般の非零Poisson点での
可換復元写像の微分や、その$7$次元核を確定するには、
17次元族のKodaira--Spencer写像、中心元の標識、
および可換復元の相対的な比較を別に確認する必要がある。
本節の観察の先行研究上の初出も未確認である。
\end{remark}
